\chapter{Vectors and Matrices}
\label{ch:matrices}

\noindent\textsc{The geometry of statistics} is linear algebra. A regression coefficient is a number, but the regression itself is a projection. A correlation is a number, but the underlying object is the cosine of an angle between vectors. Once you see this, statistical formulas stop being mysterious symbols and start being descriptions of geometric operations.

This chapter covers the linear-algebra essentials needed to read econometrics fluently.

\section{Vectors}

A vector is an ordered list of numbers:
\[
  \bm{v} = \begin{pmatrix} v_1 \\ v_2 \\ \vdots \\ v_n \end{pmatrix} \in \mathbb{R}^n.
\]
Geometrically, a vector points from the origin to a position in $n$-dimensional space.

Vectors add componentwise: $(u + v)_i = u_i + v_i$. They scale by multiplication: $(c v)_i = c \cdot v_i$.

\paragraph{The dot product.} Two vectors of the same length combine into a number:
\begin{equation}
  \bm{u} \cdot \bm{v} = \sum_{i=1}^{n} u_i v_i = u_1 v_1 + u_2 v_2 + \cdots + u_n v_n.
  \label{eq:dot-product}
\end{equation}
This is also written $\bm{u}^\top \bm{v}$ in matrix notation. The dot product is symmetric, distributive, and bilinear.

\paragraph{Length and angle.} The Euclidean length (or norm) of $\bm{v}$ is
\[
  \|\bm{v}\| = \sqrt{\bm{v} \cdot \bm{v}} = \sqrt{\sum_{i=1}^{n} v_i^2}.
\]
The angle between two vectors satisfies
\begin{equation}
  \cos\theta = \frac{\bm{u} \cdot \bm{v}}{\|\bm{u}\|\,\|\bm{v}\|}.
  \label{eq:angle-formula}
\end{equation}
Two vectors are \emph{orthogonal} (perpendicular) when their dot product is zero. Equivalently, the angle between them is $90^\circ$.

\paragraph{Why this matters for statistics.} Take vectors of mean-centered observations: $\tilde{x}_i = x_i - \bar{x}$ and $\tilde{y}_i = y_i - \bar{y}$. Then the Pearson correlation is exactly
\[
  r_{xy} = \frac{\tilde{\bm{x}} \cdot \tilde{\bm{y}}}{\|\tilde{\bm{x}}\|\,\|\tilde{\bm{y}}\|} = \cos\theta.
\]
The correlation is the cosine of the angle between two centered data vectors. $r = 1$ means the vectors point the same direction; $r = -1$ means opposite directions; $r = 0$ means orthogonal.

\section{Matrices}

A matrix is a rectangular array of numbers, organized into rows and columns:
\[
  \bm{A} = \begin{pmatrix} a_{11} & a_{12} & \cdots & a_{1n} \\ a_{21} & a_{22} & \cdots & a_{2n} \\ \vdots & \vdots & \ddots & \vdots \\ a_{m1} & a_{m2} & \cdots & a_{mn} \end{pmatrix}.
\]
This is an $m \times n$ matrix: $m$ rows, $n$ columns. We write $\bm{A} \in \mathbb{R}^{m \times n}$.

Matrices represent linear transformations. The action of $\bm{A}$ on a vector $\bm{x} \in \mathbb{R}^n$ is the matrix-vector product, a vector in $\mathbb{R}^m$:
\[
  (\bm{A}\bm{x})_i = \sum_{j=1}^{n} a_{ij} x_j.
\]
The $i$-th element of the output is the dot product of the $i$-th row of $\bm{A}$ with $\bm{x}$.

\paragraph{Matrix-matrix multiplication.} For $\bm{A} \in \mathbb{R}^{m \times n}$ and $\bm{B} \in \mathbb{R}^{n \times p}$, the product $\bm{A}\bm{B} \in \mathbb{R}^{m \times p}$ is defined by
\begin{equation}
  (\bm{A}\bm{B})_{ij} = \sum_{k=1}^{n} a_{ik} b_{kj}.
  \label{eq:matrix-product}
\end{equation}
The $(i, j)$-element is the dot product of the $i$-th row of $\bm{A}$ and the $j$-th column of $\bm{B}$. Matrix multiplication is \emph{not commutative}: in general, $\bm{A}\bm{B} \neq \bm{B}\bm{A}$.

\section{Transpose, identity, inverse}

\paragraph{Transpose.} The transpose $\bm{A}^\top$ swaps rows and columns: $(A^\top)_{ij} = A_{ji}$. For matrix products: $(\bm{A}\bm{B})^\top = \bm{B}^\top \bm{A}^\top$ (note the order reversal).

\paragraph{Identity matrix.} The matrix $\bm{I}_n$ with 1's on the diagonal and 0's elsewhere acts as a multiplicative identity: $\bm{I}_n \bm{x} = \bm{x}$ and $\bm{A}\bm{I}_n = \bm{A}$.

\paragraph{Inverse.} A square matrix $\bm{A}$ has an inverse $\bm{A}^{-1}$ if there exists $\bm{B}$ with $\bm{A}\bm{B} = \bm{B}\bm{A} = \bm{I}$. When the inverse exists, it is unique. A matrix is \emph{invertible} (or \emph{nonsingular}) iff its determinant is nonzero, iff its columns are linearly independent.

The inverse is what lets you ``solve'' matrix equations. If $\bm{A}\bm{x} = \bm{b}$ and $\bm{A}^{-1}$ exists, then $\bm{x} = \bm{A}^{-1}\bm{b}$.

\section{The OLS formula in matrix form}

This is where it all pays off. With $n$ observations and $k$ predictors:
\[
  \bm{y} = \bm{X}\bbeta + \bm{\varepsilon},
\]
where $\bm{y}$ is $n \times 1$, $\bm{X}$ is $n \times (k+1)$ (with a column of ones for the intercept), $\bbeta$ is $(k+1) \times 1$.

The OLS objective is $S(\bbeta) = \|\bm{y} - \bm{X}\bbeta\|^2 = (\bm{y} - \bm{X}\bbeta)^\top(\bm{y} - \bm{X}\bbeta)$.

Expand:
\[
  S(\bbeta) = \bm{y}^\top \bm{y} - 2\bbeta^\top \bm{X}^\top \bm{y} + \bbeta^\top \bm{X}^\top \bm{X} \bbeta.
\]
Take the gradient (vector of partial derivatives) with respect to $\bbeta$:
\[
  \nabla S(\bbeta) = -2 \bm{X}^\top \bm{y} + 2 \bm{X}^\top \bm{X} \bbeta.
\]
Set to zero:
\[
  \bm{X}^\top \bm{X} \bbeta = \bm{X}^\top \bm{y}.
\]
These are the \emph{normal equations}. If $\bm{X}^\top \bm{X}$ is invertible:
\begin{equation}
  \boxed{\hat\bbeta = (\bm{X}^\top \bm{X})^{-1} \bm{X}^\top \bm{y}.}
\end{equation}

\section{Geometry: regression as projection}

Here is the most beautiful interpretation in all of statistics.

The columns of $\bm{X}$ span a subspace of $\mathbb{R}^n$ called the \emph{column space} of $\bm{X}$, written $\text{Col}(\bm{X})$. It's the set of all vectors that can be written as $\bm{X}\bbeta$ for some choice of $\bbeta$.

The vector $\bm{y}$ generally is not in $\text{Col}(\bm{X})$. The OLS problem is: find the vector $\hat{\bm{y}} \in \text{Col}(\bm{X})$ that is closest to $\bm{y}$. The closest point is the orthogonal projection of $\bm{y}$ onto the column space.

The OLS first-order conditions are exactly the orthogonality condition: the residual $\bm{e} = \bm{y} - \hat{\bm{y}}$ is perpendicular to every column of $\bm{X}$, i.e., $\bm{X}^\top \bm{e} = \bm{0}$.

The Pythagorean theorem applies: $\|\bm{y}\|^2 = \|\hat{\bm{y}}\|^2 + \|\bm{e}\|^2$, which is the variance decomposition $SS_{\text{tot}} = SS_{\text{reg}} + SS_{\text{res}}$ when $\bm{y}$ is centered.

\begin{keyinsight}
A correlation is a cosine. A regression is a projection. The OLS estimator is the formula for orthogonal projection onto the column space of the design matrix.

Once you see this, statistical formulas become geometric statements. ``Adding a control variable'' is ``adding a dimension to the column space.'' ``$R^2$'' is ``the squared cosine of the angle between $\bm{y}$ and $\hat{\bm{y}}$.''
\end{keyinsight}

\begin{figure}[ht]
\centering
\includegraphics[width=\linewidth]{fig_vectors.pdf}
\caption{Left: vector addition. The sum $\bm{u} + \bm{v}$ is found by placing $\bm{v}$ at the tip of $\bm{u}$. Right: regression as projection. The vector $\bm{y}$ does not lie in the column space of $\bm{X}$ (the dashed line). The OLS fit $\hat{\bm{y}}$ is the orthogonal projection of $\bm{y}$ onto that subspace. The residual $\bm{e} = \bm{y} - \hat{\bm{y}}$ is perpendicular to the column space.}
\label{fig:vectors}
\end{figure}


\section{Eigenvalues and eigenvectors}

For a square matrix $\bm{A}$, an \emph{eigenvector} $\bm{v}$ and \emph{eigenvalue} $\lambda$ satisfy
\begin{equation}
  \bm{A}\bm{v} = \lambda \bm{v}.
  \label{eq:eigen}
\end{equation}
The matrix scales the eigenvector by the eigenvalue. Eigenvectors are the directions in which the transformation $\bm{A}$ acts simply.

For symmetric matrices (like covariance matrices), the eigenvalues are real and the eigenvectors form an orthonormal basis. This is the foundation of principal component analysis (PCA) and many other multivariate techniques.

\section{Practice problems}

\begin{enumerate}
  \item Compute the dot product of $\bm{u} = (1, 2, 3)$ and $\bm{v} = (4, -1, 2)$. Are they orthogonal?
  \item For $\bm{A} = \begin{pmatrix} 2 & 1 \\ 1 & 2 \end{pmatrix}$, find $\bm{A}^2$, $\bm{A}^{-1}$, and verify $\bm{A}\bm{A}^{-1} = \bm{I}$.
  \item Show that for any matrix $\bm{X}$, the matrix $\bm{X}^\top \bm{X}$ is symmetric.
  \item Set up the design matrix for the model $y = \beta_0 + \beta_1 x + \beta_2 x^2$ given three observations $(x, y) = (1, 3), (2, 7), (3, 13)$. Compute $\hat\bbeta$ using the matrix formula in R.
  \item Verify that the OLS residual vector $\bm{e}$ is orthogonal to every column of $\bm{X}$. (You should get $\bm{X}^\top \bm{e} = \bm{0}$.)
  \item Find the eigenvalues of $\bm{A} = \begin{pmatrix} 4 & 1 \\ 2 & 3 \end{pmatrix}$ by solving $\det(\bm{A} - \lambda \bm{I}) = 0$.
\end{enumerate}

\begin{reflect}
$\triangleright$~Without looking back, write the OLS formula in matrix form and explain what each piece is.

\smallskip
$\triangleright$~Why is regression a projection? What is being projected onto what?

\smallskip
$\triangleright$~The correlation between two variables equals the cosine of the angle between their centered vectors. What does this say about the relationship between $r = 0$ and orthogonality?
\end{reflect}
